%%%%%%%%%%%%%%%%%%%%%%%%%%%%%%%%%%%%%%%%%%%%%%%%%%%%%%%%%%%%%%%%%%%%%%%%%%%%%%
%  Unit II appendices
%%%%%%%%%%%%%%%%%%%%%%%%%%%%%%%%%%%%%%%%%%%%%%%%%%%%%%%%%%%%%%%%%%%%%%%%%%%%%%

\startappendices

\chapter{Syllabus-to-Chapter Concordance}

This appendix exists so that coverage can be \emph{audited} rather than assumed.
The left-hand column reproduces the Unit~II syllabus descriptor phrase by phrase,
in the order the official document prints it. The right-hand column gives the
chapter and section that carries it. Every phrase resolves to a location, and no
numbered chapter of this volume contains material that is not traceable to a
phrase in this table.

\section*{First half --- Social Media Marketing}
\addcontentsline{toc}{section}{First half --- Social Media Marketing}

\begin{mmlongtable}{Concordance: Social Media Marketing}{tab:concordanceA2}%
{@{}P{0.40\textwidth}P{0.53\textwidth}@{}}
\rowcolor{ocre}\thd{Syllabus phrase} & \thd{Where it is covered}\\
\midrule
\endfirsthead
\rowcolor{ocre}\thd{Syllabus phrase} & \thd{Where it is covered}\\
\midrule
\endhead
\rlab{Social Media Marketing: Introduction} & Chapter~1 entire. Definition and
primary goal \S\,1.1; the trust asymmetry \S\,1.2; the omnichannel gearbox
\S\,1.3; myths \S\,1.4; marketer characteristics \S\,1.5; the customer lifecycle
orbit \S\,1.6.\\
\rowcolor{tablealt}
\rlab{Major Social Media Platforms for Marketing} & Chapter~2 entire. The
deployment matrix \S\,2.1; the Facebook algorithm filter \S\,2.2; X tactics and
hashtags \S\,2.3; the LinkedIn architecture \S\,2.4; visual content \S\,2.5;
platform choice and the psychology matrix \S\,2.6.\\
\rlab{Developing Data-driven Audience \& Campaign Insights} & Chapter~3 entire.
Definition \S\,3.1; the eight-step process \S\,3.2; the persona canvas \S\,3.3;
the seven analytic layers \S\,3.4; the content lifecycle matrix \S\,3.5; the
4-1-1 rule \S\,3.6; the 4~Rs \S\,3.7.\\
\rowcolor{tablealt}
\rlab{Social Media for Business} & Chapter~4 entire. Benefits \S\,4.1; the
ten-step plan \S\,4.2; tools \S\,4.3; blogs, podcasts and webinars \S\,4.4;
monitoring \S\,4.5.\\
\rlab{Creation \& Optimization of Social Media Campaigns} & Chapter~5 entire.
Required elements \S\,5.1; the campaign build \S\,5.2; optimisation levers
\S\,5.3; metric families \S\,5.4; funnel deployment \S\,5.5.\\
\end{mmlongtable}

\section*{Second half --- Search Engine Optimization}
\addcontentsline{toc}{section}{Second half --- Search Engine Optimization}

\begin{mmlongtable}{Concordance: Search Engine Optimization}{tab:concordanceB2}%
{@{}P{0.40\textwidth}P{0.53\textwidth}@{}}
\rowcolor{ocre}\thd{Syllabus phrase} & \thd{Where it is covered}\\
\midrule
\endfirsthead
\rowcolor{ocre}\thd{Syllabus phrase} & \thd{Where it is covered}\\
\midrule
\endhead
\rlab{Search Engine Optimization Fundamentals} & Chapter~6 entire. Definition and
objective \S\,6.1; why SEO matters \S\,6.2; SEO, PPC and SEM \S\,6.3; the SERP
\S\,6.4; intent and E-E-A-T \S\,6.5; the six-step process \S\,6.6; goals and
white hat against black hat \S\,6.7.\\
\rowcolor{tablealt}
\rlab{Keywords and SEO Content Plan} & Chapter~7 entire. Why research matters
\S\,7.1; keyword types \S\,7.2; the sweet spot \S\,7.3; the ten-step content
plan and cannibalisation \S\,7.4.\\
\rlab{SEO \& Business Objectives} & Chapter~8 entire. Objective-to-KPI alignment
\S\,8.1; integrating SEO into development \S\,8.2.\\
\rowcolor{tablealt}
\rlab{Writing SEO Content} & Chapter~9 entire. Page anatomy \S\,9.1; integrating
keywords naturally \S\,9.2; key considerations \S\,9.3.\\
\rlab{On-site \& off-site SEO} & Chapter~10 entire. On-page elements \S\,10.1;
technical SEO \S\,10.2; off-page SEO \S\,10.3; the comparison \S\,10.4; the
House of SEO \S\,10.5.\\
\rowcolor{tablealt}
\rlab{Optimize Organic Search Ranking} & Chapter~11 entire. Purpose \S\,11.1;
metrics \S\,11.2; the SEO audit \S\,11.3; the human--algorithm symbiosis
\S\,11.4; tools \S\,11.5.\\
\end{mmlongtable}

\begin{keybox}[Reading the ``etc.'' in the Syllabus]
The official descriptor ends several clauses with ``etc.''. That word is treated
in this series as licence to complete a named framework, not as licence to add
new topics. Where the syllabus says ``Creation \& Optimization of Social Media
Campaigns, etc.'', the chapter covers campaign creation and optimisation
completely rather than introducing an unnamed sixth topic.
\end{keybox}

\chapter{Previous-Year Question Bank --- Unit II}

Every question in this appendix is reproduced from an actual RVCE Marketing
Management Continuous Internal Evaluation or Semester End Examination paper for
this course. Answers and pointers follow the official schemes of valuation where
those were released.

\section*{Part A --- Two-Mark Questions}
\addcontentsline{toc}{section}{Part A --- Two-Mark Questions}

\begin{enumerate}[leftmargin=2.2em,itemsep=1.1ex]

\item \textbf{What is meant by data-driven audience insight in Social Media
      Marketing?} \pyqr{CIE-II, 26 May 2026 --- 2 M}\\
      The process of collecting and analysing user data --- demographics,
      interests, online behaviour, engagement patterns and preferences --- from
      social media platforms to understand the target audience better. It helps
      marketers create personalised campaigns and improve customer engagement.
      \hfill\emph{See} \S\,3.1.

\item \textbf{List any four elements required for creating an effective social
      media campaign.} \pyqr{CIE-II, 26 May 2026 --- 2 M}\\
      Clear campaign objectives; a defined target audience; engaging and creative
      content; performance measurement and analytics. \hfill\emph{See} \S\,5.1.

\item \textbf{Define campaign optimization in Social Media Marketing and mention
      its purpose.} \pyqr{CIE-II, 26 May 2026 --- 2 M}\\
      The process of improving social media campaign performance by analysing
      metrics and adjusting content, targeting and advertisements. Purpose:
      increase engagement, improve conversion rates, reduce marketing cost and
      maximise ROI. \hfill\emph{See} \S\,5.3.

\item \textbf{List any four benefits of using social media platforms for business
      marketing and customer engagement.} \pyqr{CIE-II, 26 May 2026 --- 2 M}\\
      Increased brand awareness; better customer interaction; cost-effective
      promotion; improved customer feedback collection. \hfill\emph{See} \S\,4.1.

\item \textbf{What is the primary goal of social media marketing?}
      \pyqr{SEE, Jun/Jul 2025 --- 2 M}\\
      To build brand awareness and engagement with a target audience on social
      platforms, and to convert that engagement into traffic, leads, sales and
      long-term customer loyalty. \hfill\emph{See} \S\,1.1.

\item \textbf{Define Hashtag. List the importance of Hashtag in social media
      marketing.} \pyqr{SEE, Oct/Nov 2023 --- 2 M}\\
      A hashtag is a word or phrase prefixed with \texttt{\#} that becomes a
      clickable, searchable label grouping all posts on that topic. Importance:
      increases discoverability beyond existing followers; groups a campaign
      under one trackable label; enables participation in trending
      conversations; collects user-generated content; and provides a measurable
      unit for campaign reporting. \hfill\emph{See} \S\,2.3.

\item \textbf{Why is keyword research important in SEO? / What is keyword
      research in SEO?}
      \pyqr{CIE-II, 7 May 2025; SEE, Oct/Nov 2023 --- 2 M}\\
      Keyword research is the process of identifying the phrases that real users
      type into search engines, together with their search volume, competition
      and intent. It matters because it ensures content targets phrases that
      people actually search, that the site can realistically rank for, and that
      carry commercial intent --- so effort produces revenue rather than vanity
      traffic. \hfill\emph{See} \S\,7.1.

\item \textbf{Differentiate between On-site SEO and Off-site SEO.}
      \pyqr{CIE-II, 7 May 2025 --- 2 M}\\
      On-site SEO is work done within your own website --- title tags, meta
      descriptions, headings, content quality, keyword placement, internal
      linking, site speed --- controlling how content is presented to search
      engines. Off-site SEO is work done outside your website --- backlinks,
      digital PR, social sharing, reviews, listings --- building the site's
      authority and reputation as judged by other sites.
      \hfill\emph{See} \S\,10.4.

\item \textbf{What is the purpose of optimizing organic search rankings in SEO?}
      \pyqr{CIE-II, 7 May 2025 --- 2 M}\\
      To appear higher in unpaid results for commercially relevant queries,
      capturing sustainable high-intent traffic at near-zero marginal cost,
      reducing dependence on paid advertising and gaining the credibility of an
      organic listing. \hfill\emph{See} \S\,11.1.

\item \textbf{What is the primary objective of Search Engine Optimization?}
      \pyqr{SEE, Jun/Jul 2025 --- 2 M}\\
      To increase the quantity and quality of organic traffic by improving the
      site's visibility and ranking for relevant queries, and to ensure the site
      satisfies the intent of the visitors it attracts. \hfill\emph{See} \S\,6.1.

\item \textbf{Why is it important to integrate SEO in website development?}
      \pyqr{SEE, Jun/Jul 2025 --- 2 M}\\
      Because architecture, URLs, rendering, speed, heading structure and
      tracking are all determined at build time. Retro-fitting SEO later is
      expensive, forces URL changes that risk losing accumulated authority, and
      often leaves staging \texttt{noindex} directives in production.
      \hfill\emph{See} \S\,8.2.

\item \textbf{Name two technical aspects of a website that can significantly
      impact its organic search ranking.}
      \pyqr{CIE-II Quiz, 24 Jul 2024 --- 2 M}\\
      Site loading speed and Core Web Vitals, and mobile responsiveness. Also
      acceptable: HTTPS, crawlability and indexability, XML sitemap, URL
      structure. \hfill\emph{See} \S\,10.2.

\item \textbf{Define SEO and explain why it is important for online businesses.}
      \pyqr{CIE-II Quiz, 24 Jul 2024 --- 2 M}\\
      SEO is the set of techniques used to improve a website's visibility and
      ranking in the organic results of search engines. It matters to online
      businesses because organic search delivers about 53\% of all website
      traffic, carries the highest buyer intent of any channel, costs nothing per
      click, compounds over time, and confers more credibility than paid
      listings. \hfill\emph{See} \S\,6.1 and \S\,6.2.

\item \textbf{What are two key considerations when writing SEO-friendly content
      for web pages?} \pyqr{CIE-II Quiz, 24 Jul 2024 --- 2 M}\\
      (1) Match the search intent --- the format and depth must answer what the
      user actually wants. (2) Integrate keywords naturally in the title, H1,
      opening paragraph and subheadings using semantic variants, at roughly
      3--5\% density, without stuffing. Also acceptable: a unique title tag and
      meta description; a scannable structure. \hfill\emph{See} \S\,9.3.

\end{enumerate}

\section*{Part B --- Eight- and Ten-Mark Questions}
\addcontentsline{toc}{section}{Part B --- Eight- and Ten-Mark Questions}

\begin{enumerate}[leftmargin=2.2em,itemsep=1.4ex]

\item \textbf{Analyse how the Facebook Algorithm Filter evaluates content
      distribution. What specific practices are penalized (filtered out) and
      which ones are rewarded (propelled through)?}
      \pyqr{CIE-II, 26 May 2026 --- 10 M, both sets}\\
      Structure the answer as: (i) what the algorithm is and why it exists;
      (ii) the five evaluation criteria --- engagement history, relevance,
      content type preference, timeliness, quality and originality; (iii) the
      penalised list --- clickbait headlines, repeated sources posted
      back-to-back, engagement bait, fake news, spam, low-quality copied
      content, excessive promotion; (iv) the rewarded list --- time spent on the
      post after full load, video actions (unmuting, full-screen, HD), Facebook
      Live at 300\% more watch time than recorded video, original content,
      meaningful interactions, informative posts, consistent posting.
      \hfill\emph{See} \S\,2.2.

\item \textbf{Based on ``The Content Lifecycle Matrix,'' how should a marketer
      adjust the objective, tone, and format of their content as a prospect
      transitions from the Early-Stage to the Late-Stage?}
      \pyqr{CIE-II, 26 May 2026 --- 10 M}\\
      Reproduce the three-by-three matrix and then explain the trajectory: the
      objective moves from attention to trust to action; the tone moves from fun
      and entertaining to educational to informative and conversion-focused; the
      format moves from curated news, tips and visual content to contests,
      newsletter subscriptions and event invitations, and finally to
      click-to-purchase, demonstrations and trial sign-ups.
      \hfill\emph{See} \S\,3.5.

\item \textbf{Explain the stages of ``The Customer Lifecycle Orbit'' and discuss
      the statistical differences in selling to new prospects versus existing
      customers.} \pyqr{CIE-II, 26 May 2026 --- 10 M}\\
      Name and explain the six stages, then give the statistics: the probability
      of selling to a new prospect is roughly 5--20\%, against 60--70\% for an
      existing customer; repeat customers spend up to 67\% more; 49\% of
      marketers achieve higher ROI by focusing on engagement. Conclude with the
      implication --- retention is structurally cheaper and more profitable than
      acquisition, so social content must serve the whole orbit, not just
      awareness. \hfill\emph{See} \S\,1.6.

\item \textbf{Detail the ``4 Rs of Content Scaling.'' How does this strategic
      framework help marketers maximize the use of their content assets?}
      \pyqr{CIE-II, 26 May 2026 --- 10 M}\\
      Give Reorganize, Rewrite, Retire and Redesign with an example each, note
      the alternative Repurpose, Refresh, Repost and Recombine formulation, and
      close with the benefits --- saves time and cost, increases reach, improves
      SEO performance, enhances engagement and maximises ROI from content
      assets. \hfill\emph{See} \S\,3.7.

\item \textbf{Explain the importance of social media marketing for businesses
      today. How can businesses integrate social media platforms into their
      overall marketing strategy?} \pyqr{CIE-II, 7 May 2025 --- 10 M}\\
      \emph{Importance} --- the trust asymmetry (90\% believe peers, 48\% trust
      businesses), reach, cost-effectiveness, two-way engagement, feedback,
      social selling and SEO fuel. \emph{Integration} --- the omnichannel
      gearbox, the ten-step plan, mapping content to the customer lifecycle
      orbit, and synchronising social with website, e-mail, search and offline
      channels so the customer experiences one continuous conversation.
      \hfill\emph{See} \S\,1.2, \S\,1.3 and \S\,4.2.

\item \textbf{Compare the roles of Facebook, Instagram, LinkedIn, and Twitter in
      social media marketing. How should a business choose the right platform?}
      \pyqr{CIE-II, 7 May 2025 --- 10 M}\\
      Use the platform deployment matrix covering core mechanic, best use and key
      features for each, then the six selection criteria, closing with the
      platform psychology matrix --- each platform serves a different
      psychological function in the decision. \hfill\emph{See} \S\,2.1 and
      \S\,2.6.

\item \textbf{Explain the importance of data-driven insights in social media
      marketing. Describe the process of developing audience insights using
      social media analytics tools. How can these insights inform the
      development and optimization of social media campaigns?}
      \pyqr{CIE-I, 1 Jul 2024 --- 10 M}\\
      Cover the definition, the eight-step insight process, the tools, and then
      the application --- targeting, creative, tone, format, schedule and budget
      allocation, plus in-flight optimisation using the eight levers.
      \hfill\emph{See} Chapter~3 and \S\,5.3.

\item \textbf{Illustrate how social media analytics contribute to creating
      audience profiles.} \pyqr{SEE, Jun/Jul 2025, Q4a --- 8 M}\\
      The seven layers: demographic, temporal, interest, behavioural,
      conversational, comparative and conversion.
      \hfill\emph{See} \S\,3.4.

\item \textbf{Analyse how businesses can utilize Facebook and Instagram for
      different stages of a marketing funnel.}
      \pyqr{SEE, Jun/Jul 2025, Q3a --- 8 M}\\
      Cover awareness, consideration, conversion and retention or advocacy on
      both platforms. \hfill\emph{See} \S\,5.5.

\item \textbf{Discuss the key metrics used to measure the success of a social
      media marketing campaign and explain the importance of it for evaluating
      performance.} \pyqr{SEE, Oct/Nov 2023, Q3a --- 8 M}\\
      The six metric families and the six reasons measurement matters.
      \hfill\emph{See} \S\,5.4.

\item \textbf{Analyse the differences between on-site and off-site SEO
      techniques. Provide examples for each.}
      \pyqr{CIE-II, 7 May 2025 --- 10 M}\\
      Use the comparison table and the House of SEO, with the worked examples
      given there. \hfill\emph{See} \S\,10.4 and \S\,10.5.

\item \textbf{Explain the key components of Search Engine Optimization,
      including on-page, off-page, and technical SEO. How do these elements
      contribute to improving website ranking on search engines?}
      \pyqr{CIE-I, 16 Apr 2026 --- 10 M}\\
      \emph{On-page SEO} (content optimisation) --- keyword optimisation in
      titles, headings and content; meta tags; URL structure; internal linking;
      high-quality relevant content. \emph{Off-page SEO} (external factors) ---
      backlinks from other websites; social media sharing; influencer
      collaborations; online reviews. \emph{Technical SEO} (backend
      optimisation) --- website speed; mobile friendliness; secure HTTPS; XML
      sitemap; crawlability and indexing. Then explain the contribution:
      technical SEO makes the site discoverable, on-page SEO makes it relevant,
      and off-page SEO makes it authoritative --- and rankings are a function of
      all three together. \hfill\emph{See} Chapter~10.

\item \textbf{What are the key metrics used to measure organic search ranking?
      How can businesses monitor and improve these metrics over time? Describe
      the process of conducting an SEO audit and implementing changes to optimize
      organic search ranking.} \pyqr{CIE-II, 24 Jul 2024 --- 10 M}\\
      The metric table, then the nine-step audit process, closing with the
      six-step SEO process as the ongoing improvement cycle.
      \hfill\emph{See} \S\,11.2, \S\,11.3 and \S\,6.6.

\item \textbf{An educational institution is struggling to rank for competitive
      keywords related to online courses. Conduct an SEO audit of their website,
      identify key issues affecting organic search ranking, and recommend
      actionable solutions.} \pyqr{CIE-II, 24 Jul 2024 --- 10 M}\\
      The worked diagnostic is set out in full in the chapter.
      \hfill\emph{See} \S\,11.3.

\item \textbf{You are writing a blog post about a complex technical topic.
      Describe how you would integrate relevant keywords naturally into the
      content while ensuring it remains engaging and informative for readers.}
      \pyqr{CIE-II, 24 Jul 2024 --- 10 M}\\
      The nine-point method, anchored on ``write for the reader first, then
      optimise'': keyword placement in high-value positions, semantic variants,
      3--5\% density, People-Also-Ask coverage, jargon explanation, visual
      breaks, internal linking and a read-aloud check.
      \hfill\emph{See} \S\,9.2.

\item \textbf{Develop an SEO content plan for an e-commerce startup targeting
      eco-friendly household products.} \pyqr{SEE, Jun/Jul 2025, Q3b --- 8 M}\\
      \emph{Objective:} organic revenue from high-intent shoppers within nine
      months. \emph{Keyword research:} cluster into topics --- eco-friendly
      cleaning, plastic-free kitchen, sustainable laundry, bamboo and
      natural-fibre goods --- targeting long-tail commercial phrases such as
      ``biodegradable dishwashing liquid India'' (volume 250--1,000+, difficulty
      under 40, clear buyer intent). \emph{Architecture:} one pillar page per
      cluster, supported by product and comparison pages, all internally linked.
      \emph{Content by intent:} informational blogs (``how harmful are
      conventional detergents?''), commercial comparisons (``best plastic-free
      kitchen swaps''), and transactional category and product pages with unique
      copy. \emph{On-page:} unique title tags under 70 characters, 155--160
      character meta descriptions, one H1 per page, ALT text on all product
      images, 3--5\% keyword density, breadcrumb navigation.
      \emph{Technical:} a fast mobile-first storefront, HTTPS, XML sitemap,
      Product and Review schema for rich results, canonical tags on filtered
      category URLs. \emph{Off-page:} digital PR on plastic-waste data, guest
      posts on sustainability blogs, Google Business Profile, review generation.
      \emph{Measurement:} keyword coverage, organic sessions, organic conversion
      rate and revenue, reviewed monthly. \hfill\emph{See} \S\,7.4.

\item \textbf{Construct a targeted marketing plan using SEO tools like Google
      Analytics and SEMrush for a B2B company.}
      \pyqr{SEE, Jun/Jul 2025, Q4b --- 8 M}\\
      \emph{Step 1 --- baseline in Google Analytics 4:} identify current organic
      sessions, top landing pages, acquisition channels, and which pages assist
      conversions. \emph{Step 2 --- opportunity analysis in SEMrush:} run
      Organic Research on the top five competitors, use Keyword Gap to find
      phrases they rank for and you do not, and the Keyword Magic Tool to build
      clusters around the buyer's problems. \emph{Step 3 --- persona and intent
      mapping:} align keywords to the four to six B2B personas (Executive
      Sponsor, Decision Maker, End User), matching informational content to
      early-stage roles and commercial content to decision makers.
      \emph{Step 4 --- content plan:} pillar pages per solution area, supported
      by case studies, ROI calculators, whitepapers and comparison pages.
      \emph{Step 5 --- technical and on-page fixes} from the SEMrush Site Audit,
      prioritised by impact. \emph{Step 6 --- off-page:} industry publication
      guest posts, original research for digital PR, LinkedIn thought
      leadership. \emph{Step 7 --- measurement:} SEMrush Position Tracking for
      rankings and share of voice; GA4 goals and UTM-tagged campaigns for
      marketing-qualified leads, cost per lead and pipeline influence; monthly
      review and reallocation. \hfill\emph{See} \S\,11.5 and \S\,8.1.

\item \textbf{Explain the concept of Search Engine optimization and its
      significance in digital marketing.}
      \pyqr{SEE, Oct/Nov 2023, Q4a --- 8 M}\\
      Define SEO, explain the SEO, PPC and SEM distinction, the six-step process,
      and the three pillars --- technical, on-page and off-page.
      \emph{Significance:} 53\% of all traffic, the highest buyer intent,
      compounding at near-zero marginal cost, greater credibility than
      advertisements, competitive defence, and measurable results.
      \hfill\emph{See} Chapter~6.

\item \textbf{Describe the role of content in SEO. How can business create high
      quality SEO friendly content?}
      \pyqr{SEE, Oct/Nov 2023, Q4b --- 8 M}\\
      \emph{Role:} content is king --- search engines store and rank text, so
      high-volume, high-quality relevant content both gives users a reason to
      stay and return, and gives engines the material from which rankings are
      computed. It is also the asset that earns backlinks and satisfies E-E-A-T.
      \emph{How:} research keywords and intent first; write for the reader and
      optimise afterwards; use the anatomy of a top-ranking page; maintain
      3--5\% keyword density with semantic variants; write unique title tags and
      meta descriptions; structure with one H1 and logical H2 and H3 headings;
      add rich media for dwell time; link internally with descriptive anchors;
      keep content fresh through historical audits; and ensure it is more
      comprehensive and more readable than what currently ranks.
      \hfill\emph{See} Chapter~9.

\item \textbf{Enumerate briefly on the best practices for creating compelling and
      shareable content on social media platforms to increase brand visibility
      and engagement.} \pyqr{SEE, Oct/Nov 2023, Q3b --- 8 M}\\
      Apply the seven-point quality gate from Volume~1; lead with visuals ---
      40 times more shareable, 94\% more views with relevant images, and 135\%
      greater organic reach for video; follow the 4-1-1 rule; write natively for
      each platform rather than cross-posting; hook within the first three
      seconds; use one or two hashtags only; post consistently at
      audience-active times; invite participation through polls, questions and
      contests; scale one master asset with the 4~Rs; and amplify through
      communities, employees and influencers.
      \hfill\emph{See} \S\,2.5, \S\,3.6 and \S\,3.7.

\end{enumerate}

\section*{How Unit II Is Examined}
\addcontentsline{toc}{section}{How Unit II Is Examined}

\begin{keybox}[The Pattern Across Papers]
In the Semester End Examination, Unit~II is Questions~3 and~4 with internal
choice --- 16 marks in two eight-mark sub-divisions --- plus two or three Part~A
objective questions. The second Continuous Internal Evaluation test is typically
dominated by this unit.

\smallskip
The highest-probability ten-mark topics are the \textbf{Facebook algorithm
filter} (asked twice in one paper), the \textbf{content lifecycle matrix}, the
\textbf{4~Rs of content scaling}, the \textbf{customer lifecycle orbit} with its
statistics, \textbf{on-site against off-site SEO}, and the \textbf{SEO audit
process}.

\smallskip
Note that the Semester End Examination also mixes social media marketing and SEO
into applied ``design a plan for X'' questions. Always answer those with a named
framework, not free-form opinion.
\end{keybox}

\chapter{Supplementary Topic: Search Everywhere Optimization}

\begin{warnbox}[What This Appendix Is, and Why It Is Not a Chapter]
Search Everywhere Optimization --- sometimes called SEO 2.0 --- appears in the
supplementary lecture material for this unit but is named by no unit descriptor
in the official syllabus. It is placed in an appendix so that the rule governing
this series holds: the numbered chapters carry the syllabus and nothing else. It
is included rather than dropped because it is directly useful for any question
that asks about \emph{new} digital channels or about why traffic can rise while
conversions stay flat.
\end{warnbox}

\section{The Scale of Search Outside Google}

\begin{keybox}[The Figures]
Around \textbf{13.7 billion} searches happen daily across the internet, and
\textbf{73\%} of all searches happen \emph{outside} Google. Most marketers are
still playing a Google-only game --- which means a business winning at Google can
simultaneously be losing customers.
\end{keybox}

\section{The Symptom and the Disease}

\begin{description}
\item[The symptom] Traffic looks decent but conversions are flat. The business is
showing up for the search and missing the decision.
\item[The disease] Optimising for visibility in one place while the customer is
deciding everywhere else.
\end{description}

\section{The Death of the Funnel}

The modern journey is not funnel-shaped and linear. It is a \emph{decision
constellation} of micro-choices happening in different places at the same time.

\begin{figure}[htbp]
\centering
\begin{tikzpicture}[font=\sffamily\footnotesize,
    nd/.style={draw=ocre!70,line width=0.8pt,fill=ocrepale,rounded corners=3pt,
               align=center,text width=2.05cm,inner sep=4pt,
               font=\sffamily\scriptsize,text=slate}]
  \node[draw=slate,line width=1.1pt,fill=white,circle,minimum size=2.3cm,
        align=center,font=\sffamily\bfseries\scriptsize,text=slate] (c) at (0,0)
        {THE\\ DECISION};
  \foreach \ang/\n/\lbl in {%
      90/a/{\textbf{Click}\\ search engine},
      45/b/{\textbf{Trust}\\ forums},
      0/c2/{\textbf{Buy}\\ marketplaces},
     -45/d/{\textbf{Try}\\ app stores},
     -90/e/{\textbf{Think}\\ video, podcasts},
    -135/f/{\textbf{Believe}\\ AI assistants},
     180/g/{\textbf{Follow}\\ social feeds},
     135/h/{\textbf{Better}\\ video reviews}}{
     \node[nd] (\n) at (\ang:3.75cm) {\lbl};
     \draw[ocre!70,line width=0.7pt,dashed] (\n) -- (c);}
\end{tikzpicture}
\caption[The decision constellation]%
        {The decision constellation. Because these micro-decisions happen in
         parallel rather than in sequence, optimising a single channel optimises
         a single point in a distributed process.}
\label{fig:constellation}
\end{figure}

\section{Visibility Against Validation}

\begin{keybox}[The Distinction]
\textbf{Visibility} is what you do --- showing up in the search engine, creating
platform-specific posts. \textbf{Validation} is what others say about what you do
--- organic mentions in conversations, authentic social proof.

\smallskip
\textbf{Visibility without validation is just noise.} You need both to secure the
decision.
\end{keybox}

\section{Prioritisation: the RICE Framework}

You do not have to be everywhere. You need a prioritisation strategy.

\begin{mmlongtable}{The RICE prioritisation framework}{tab:rice}%
{@{}P{0.16\textwidth}P{0.77\textwidth}@{}}
\rowcolor{ocre}\thd{Criterion} & \thd{The question it asks}\\
\midrule
\endfirsthead
\rowcolor{ocre}\thd{Criterion} & \thd{The question it asks}\\
\midrule
\endhead
\rlab{Reach} & How many users in your specific audience will this
micro-decision impact?\\
\rowcolor{tablealt}
\rlab{Impact} & How deeply does this platform influence the final purchase
decision?\\
\rlab{Confidence} & How certain are you that your brand can execute
authentically here?\\
\rowcolor{tablealt}
\rlab{Ease} & What operational effort is required to maintain a strategic
presence?\\
\end{mmlongtable}

\chapter{Glossary of Key Terms}

\begin{mmlongtable}{Key terms for rapid revision}{tab:glossary2}%
{@{}P{0.27\textwidth}P{0.66\textwidth}@{}}
\rowcolor{ocre}\thd{Term} & \thd{One-line meaning}\\
\midrule
\endfirsthead
\rowcolor{ocre}\thd{Term} & \thd{One-line meaning}\\
\midrule
\endhead
\rlab{4-1-1 rule} & Four value posts, one soft promotion, one hard promotion
(Joe Pulizzi)\\
\rowcolor{tablealt}
\rlab{4 Rs of content scaling} & Reorganize, Rewrite, Retire, Redesign\\
\rlab{Backlink} & An inbound link from another site --- the currency of off-page
SEO\\
\rowcolor{tablealt}
\rlab{Campaign optimization} & Improving campaign performance by analysing
metrics and adjusting content, targeting and advertisements\\
\rlab{Content lifecycle matrix} & Early, mid and late-stage mapping of objective,
tone and format\\
\rowcolor{tablealt}
\rlab{Core Web Vitals} & LCP, INP and CLS --- the technical experience signals
that feed into ranking\\
\rlab{Customer lifecycle orbit} & Awareness, engagement, purchase, retention and
loyalty, growth, advocacy\\
\rowcolor{tablealt}
\rlab{Data-driven audience insight} & Collecting and analysing platform user data
to understand the audience\\
\rlab{E-E-A-T} & Experience, Expertise, Authoritativeness, Trustworthiness\\
\rowcolor{tablealt}
\rlab{Engagement bait} & Explicitly asking for likes, shares or comments ---
penalised by the algorithm\\
\rlab{Facebook algorithm filter} & The ranking system deciding which posts reach
the news feed\\
\rowcolor{tablealt}
\rlab{Featured snippet} & ``Position zero'' --- a direct answer extracted onto
the SERP\\
\rlab{Hashtag} & A \texttt{\#}-prefixed searchable label grouping posts on a
topic\\
\rowcolor{tablealt}
\rlab{Head keyword} & A one- or two-word phrase: huge volume, brutal
competition, vague intent\\
\rlab{Keyword cannibalisation} & Two of your own pages competing for the same
keyword\\
\rowcolor{tablealt}
\rlab{Keyword density} & Share of the content occupied by the focus keyword;
industry norm 3--5\%\\
\rlab{Link gravity} & The principle that link quality and relevance outweigh raw
link count\\
\rowcolor{tablealt}
\rlab{Long-tail keyword} & A four-or-more-word phrase: low volume, low
competition, high intent\\
\rlab{Off-page SEO} & A site's authority as determined by what other sites say
about it\\
\rowcolor{tablealt}
\rlab{On-page SEO} & How well a site's own content and code are presented to
search engines\\
\rlab{Panda / Penguin} & Google updates penalising thin content and manipulative
link-building\\
\rowcolor{tablealt}
\rlab{RankBrain / Hummingbird} & Google algorithms, from 2015, focused on user
intent rather than keyword strings\\
\rlab{RICE} & Reach, Impact, Confidence, Ease --- a channel prioritisation
framework\\
\rowcolor{tablealt}
\rlab{SEM} & Umbrella term covering SEO plus PPC\\
\rlab{SEO} & Improving organic visibility and ranking in search engines\\
\rowcolor{tablealt}
\rlab{SERP} & Search engine results page\\
\rlab{Share of voice} & Conversations about your brand against conversations
about competitors\\
\rowcolor{tablealt}
\rlab{Social media marketing} & Creating and sharing content on social networks
to achieve marketing and branding goals\\
\rlab{Social media monitoring} & Tracking mentions of brand, competitors,
keywords and hashtags\\
\rowcolor{tablealt}
\rlab{Technical SEO} & The crawlable, indexable, fast, secure infrastructure of
the site\\
\rlab{Trust asymmetry} & Roughly 90\% believe peers against about 48\% who trust
businesses directly\\
\rowcolor{tablealt}
\rlab{White hat / black hat SEO} & Techniques within search-engine guidelines,
against techniques that violate them\\
\end{mmlongtable}

\chapter{Framework Quick-Reference Sheet}

Everything in this volume that is a numbered list, in one place, for the night
before the examination.

\begin{mmlongtable}{Framework quick reference}{tab:quickref2}%
{@{}P{0.28\textwidth}P{0.65\textwidth}@{}}
\rowcolor{ocre}\thd{Framework} & \thd{The sequence}\\
\midrule
\endfirsthead
\rowcolor{ocre}\thd{Framework} & \thd{The sequence}\\
\midrule
\endhead
\rlab{Trust asymmetry} & About 90\% believe peers; about 48\% trust businesses;
word of mouth valued about 41\% above social recommendations\\
\rowcolor{tablealt}
\rlab{Omnichannel gearbox} & Social selling, SEO fuel, omnichannel
synchronisation\\
\rlab{Customer lifecycle orbit} & Awareness \ra{} engagement \ra{} purchase
\ra{} retention and loyalty \ra{} growth \ra{} advocacy\\
\rowcolor{tablealt}
\rlab{Retention statistics} & New prospect under 20\% (5--20\%); existing
customer over 60\% (60--70\%); repeat customers spend up to 67\% more; 49\%
report higher ROI from engagement\\
\rlab{Eight-step insight process} & Define the question \ra{} collect \ra{}
consolidate \ra{} segment \ra{} find patterns \ra{} build personas \ra{} apply
\ra{} iterate\\
\rowcolor{tablealt}
\rlab{Seven analytic layers} & Demographic, temporal, interest, behavioural,
conversational, comparative, conversion\\
\rlab{Persona canvas} & Background and goals; information sources; messaging and
topics; purchase role and objections\\
\rowcolor{tablealt}
\rlab{Content lifecycle matrix} & Objective: attention \ra{} trust \ra{} action.
Tone: fun \ra{} educational \ra{} decisive. Format: curated \ra{} participatory
\ra{} transactional\\
\rlab{4-1-1 rule} & Four value + one soft promotion + one hard promotion per six
posts\\
\rowcolor{tablealt}
\rlab{4 Rs of content scaling} & Reorganize, Rewrite, Retire, Redesign --- also
described as Repurpose, Refresh, Repost, Recombine\\
\rlab{Facebook filter --- rewarded} & Time spent, video actions, Live, original
content, meaningful interaction, informative posts, consistency\\
\rowcolor{tablealt}
\rlab{Facebook filter --- penalised} & Clickbait, repeated sources, engagement
bait, fake news, spam, copied content, excessive promotion\\
\rlab{Platform choice criteria} & Persona location, B2C against B2B, format
capability, funnel stage, resources, measured ROI\\
\rowcolor{tablealt}
\rlab{Platform psychology} & Discovery, education, verification, purchasing,
aspiration, validation, the click\\
\rlab{Ten-step SMM plan} & Goals \ra{} audit \ra{} audience \ra{} platforms
\ra{} content strategy \ra{} calendar \ra{} resources \ra{} publish and engage
\ra{} amplify \ra{} measure\\
\rowcolor{tablealt}
\rlab{Campaign elements} & Clear objectives, defined audience, engaging creative,
measurement\\
\rlab{Eight optimisation levers} & Creative, copy and CTA, audience, placement,
timing and frequency, budget, landing page, format\\
\rowcolor{tablealt}
\rlab{Six metric families} & Reach and volume, engagement, traffic, conversion,
sentiment and advocacy, efficiency\\
\rlab{SEO, PPC, SEM} & Organic, paid, and the umbrella over both\\
\rowcolor{tablealt}
\rlab{SERP result types} & Paid, organic, universal or vertical\\
\rlab{Six-step SEO process} & Keyword research \ra{} reporting and goals \ra{}
content building \ra{} page optimisation \ra{} social and link building \ra{}
follow-up analysis\\
\rowcolor{tablealt}
\rlab{Keyword intent types} & Informational, navigational, commercial
investigation, transactional\\
\rlab{Keyword lengths} & Head (1--2), body (2--3), long tail (4+)\\
\rowcolor{tablealt}
\rlab{Keyword sweet spot} & Volume 250--1,000+ $\times$ difficulty under 40
$\times$ commercial intent\\
\rlab{Ten-step content plan} & Objectives \ra{} keyword clusters \ra{} intent
mapping \ra{} content audit \ra{} gap analysis \ra{} pillar architecture \ra{}
one keyword per page \ra{} calendar \ra{} targets \ra{} monthly review\\
\rowcolor{tablealt}
\rlab{On-page five elements} & Page copy, title tags, meta data, heading tags,
interlinking\\
\rlab{On-page limits} & Title about 70 characters; meta description 155--160
characters; keyword density 3--5\%; one H1 per page\\
\rowcolor{tablealt}
\rlab{House of SEO} & Technical (foundation, discoverable), on-page (floor,
relevant), off-page (roof, authoritative)\\
\rlab{Nine-step SEO audit} & Crawl \ra{} indexation \ra{} technical \ra{}
on-page \ra{} content \ra{} backlinks \ra{} benchmark \ra{} prioritise \ra{}
implement and re-measure\\
\rowcolor{tablealt}
\rlab{RICE (Appendix~C)} & Reach, Impact, Confidence, Ease\\
\end{mmlongtable}
